\documentclass[11pt]{article}

\usepackage[margin=1in]{geometry}
\usepackage{amsmath,amssymb,amsfonts,amsthm}
\usepackage{array}
\IfFileExists{bbm.sty}{\usepackage{bbm}}{}
\usepackage{bm}
\usepackage{graphicx}
\usepackage{multirow}
\usepackage{multicol}
\usepackage{nicefrac}
\usepackage{paralist}
\usepackage{enumitem}
\usepackage{verbatim}
\usepackage{thm-restate}
\usepackage[normalem]{ulem}
\usepackage[compact]{titlesec}
\usepackage{tikz}
\usetikzlibrary{arrows}
\usetikzlibrary{arrows.meta}
\usetikzlibrary{backgrounds}
\usetikzlibrary{calc}
\usetikzlibrary{decorations.markings}
\usetikzlibrary{fit}
\usetikzlibrary{patterns}
\usetikzlibrary{positioning}
\usetikzlibrary{shapes}

\usepackage{fancybox}
\newenvironment{fminipage}%
  {\begin{Sbox}\begin{minipage}}%
  {\end{minipage}\end{Sbox}\fbox{\TheSbox}}

\newenvironment{algbox}[0]{%
\vskip 0.2in
\noindent
\begin{fminipage}{6.3in}
}{%
\end{fminipage}
\vskip 0.2in
}

\IfFileExists{algorithm2e.sty}{\usepackage[ruled]{algorithm2e}}{}
\IfFileExists{algpseudocode.sty}{\usepackage[noend]{algpseudocode}}{}
\usepackage{tcolorbox}
\tcbuselibrary{skins,breakable}
\tcbset{enhanced jigsaw}
\usepackage{mdframed}
\usepackage{xcolor}
\usepackage{cite}
\usepackage{url}
\usepackage{xspace}
\usepackage[mathscr]{euscript}
\usepackage{nameref}
\usepackage[
  linktocpage=true,
  pagebackref=true,
  colorlinks=true,
  linkcolor=blue,
  citecolor=blue,
  urlcolor=blue,
  bookmarks,
  bookmarksopen,
  bookmarksnumbered
]{hyperref}
\usepackage[noabbrev,nameinlink]{cleveref}

\newtheorem{theorem}{Theorem}[section]
\newtheorem{example}{Example}[section]
\newtheorem{lemma}{Lemma}[section]
\newtheorem{proposition}[lemma]{Proposition}
\newtheorem{corollary}[lemma]{Corollary}
\newtheorem{claim}[lemma]{Claim}
\newtheorem{fact}[lemma]{Fact}
\newtheorem{conjecture}[lemma]{Conjecture}
\newtheorem{definition}[lemma]{Definition}
\newtheorem{problem}{Problem}
\newtheorem{remark}[lemma]{Remark}
\newtheorem{observation}[lemma]{Observation}

\newtheorem*{theorem*}{Theorem}
\newtheorem*{claim*}{Claim}
\newtheorem*{proposition*}{Proposition}
\newtheorem*{lemma*}{Lemma}
\newtheorem*{problem*}{Problem}

\allowdisplaybreaks
\setlength{\parskip}{3pt}

\title{Regular Distributions Are Not Closed Under Convolution}
\author{}
\date{August 17, 2026}

\begin{document}

\pagenumbering{roman}

\maketitle

\begin{abstract}
We show that Myerson regularity is not preserved by convolution, even for two
independent identically distributed bounded nonnegative summands.  For an
absolutely continuous distribution $F$ with density $f$, regularity means that
the virtual value $x-(1-F(x))/f(x)$ is nondecreasing at points where $f$ and the
survival probability are positive.  Let
$a=1/50$, $c=49/100$, $d=51/100$, and $B=a/d^2=200/2601$.  We define a common
marginal law supported on $[0,1]$ by the survival function
\[
\overline F(x)=
\begin{cases}
\dfrac{a}{a+x}, & 0 \leq x \leq c,\\[6pt]
B(1-x), & c \leq x \leq 1,
\end{cases}
\]
with $\overline F(x)=1$ for $x<0$ and $\overline F(x)=0$ for $x>1$.  This
marginal is regular, with virtual value $-a$ on $[0,c]$ and $2x-1$ on
$[c,1)$.  If $S$ is the sum of two independent copies and $\Phi$ is the virtual
value of $S$, then
\[
\left.\frac{d}{dt}\Phi(1+t)\right|_{t=0+}
=
\frac{31255005}{12510002}
-\frac{6759999}{6255001}\log\frac{51}{2}
<0.
\]
Thus $\Phi$ decreases immediately to the right of the interior point $1$ of the
support of $S$, and the distribution of $S$ is not regular.
\end{abstract}

\noindent\textbf{Note.} This paper was fully AI generated by GPT 5.5 at xhigh reasoning effort using an internal harness, prompted by Richard Peng.\par



\clearpage
\tableofcontents
\clearpage

\pagenumbering{arabic}

\section{Introduction}
\label{sec:introduction}

Regularity is one of the basic shape assumptions in optimal auction theory.
For an absolutely continuous distribution $F$ with density $f$ and survival
function $\overline F=1-F$, the associated virtual value is
\[
\varphi_F(x)=x-\frac{\overline F(x)}{f(x)}
\]
at points where both $f(x)$ and $\overline F(x)$ are positive.  The
distribution is regular, in the sense of Myerson, if $\varphi_F$ is
nondecreasing on the part of the support on which it is defined
\cite{Myerson1981OptimalAuction}.  This
monotonicity is what lets optimal-auction rules be expressed without ironing;
the same virtual-value calculus also appears in the classical marginal-revenue
interpretation of optimal auctions and monopoly pricing
\cite{RileySamuelson1981OptimalAuctions,BulowRoberts1989SimpleEconomics}.

It is natural to ask whether regularity is preserved by aggregation.  Given
independent nonnegative values $X_1,\ldots,X_m$, each with a regular marginal,
does the sum
\[
S=X_1+\cdots+X_m
\]
again have a regular distribution?  Equivalently, is the class of
Myerson-regular distributions closed under convolution?  This is not a formal
consequence of the marginal definition: the virtual value of the sum depends on
the ratio of the convolution survival function to the convolution density, and
that ratio need not inherit monotonicity from the corresponding marginal
ratios.

Positive closure results are known for stronger or different distributional
hypotheses.  The closest benchmark for the same aggregation operation is the
monotone-hazard-rate, or IFR/MHR, condition: if the hazard rate
$h_F=f/\overline F$ is nondecreasing, then $1/h_F$ is nonincreasing and
$x-1/h_F(x)$ is nondecreasing, so MHR implies Myerson regularity.  Since
IFR/MHR is a stronger assumption placed on the same convolution question,
Barlow, Marshall, and Proschan's closure theorem for IFR/MHR distributions is
the main positive restricted comparison \cite{BarlowMarshallProschan1963MHR}.
Log-concavity is another common positive regime: the Pr{\'e}kopa theorem gives
stability under marginalization and convolution, and standard one-dimensional
log-concavity facts imply stronger survival and hazard-rate properties in
economic applications
\cite{Prekopa1973LogConcave,BagnoliBergstrom2005LogConcave}.  Separately,
regularity and its quantitative strengthenings have been studied in mechanism
design under formulations such as generalized concavity, $\alpha$-strong
regularity, and related performance assumptions
\cite{Ewerhart2013RegularType,ColeRoughgarden2014SampleComplexity,ColeRao2015AlphaStronglyRegular,SchweizerSzech2019PerformanceBounds}.
Revenue-management assumptions such as increasing generalized failure rate are
nearby but different, since they involve monotonicity of $x h_F(x)$ and related
pricing quantities rather than monotonicity of $x-1/h_F(x)$
\cite{Lariviere2006IGFR,ZiyaAyhanFoley2004RevenueManagement,BanciuMirchandani2013IGFR}.
These works provide important context and stronger or adjacent sufficient
conditions, but they do not imply convolution closure for ordinary Myerson
regularity.

The contribution of this paper is an explicit counterexample showing that the
ordinary regularity assumption is too weak for such a closure theorem.  The
failure already occurs for the smallest nontrivial convolution, with two
independent identically distributed summands supported on $[0,1]$.

\begin{theorem}[Explicit i.i.d. convolution counterexample]
\label{thm:main}
Let
\[
a=\frac1{50},\qquad c=\frac{49}{100},\qquad d=\frac{51}{100},
\qquad B=\frac{a}{d^2}=\frac{200}{2601}.
\]
There exist independent identically distributed nonnegative continuously
distributed random variables $X_1,X_2$, supported on $[0,1]$, whose common
survival function on $0\leq x\leq 1$ is
\[
\overline F(x)=
\begin{cases}
\dfrac{a}{a+x}, & 0 \leq x \leq c,\\[6pt]
B(1-x), & c \leq x \leq 1
\end{cases}
\]
with $\overline F(x)=1$ for $x<0$ and $\overline F(x)=0$ for $x>1$.
The common distribution is regular.  However, if $S=X_1+X_2$ and $\Phi$ is the
virtual value of the distribution of $S$, then
\[
\left.\frac{d}{dt}\Phi(1+t)\right|_{t=0+}
=
\frac{31255005}{12510002}
-\frac{6759999}{6255001}\log\frac{51}{2}
<0.
\]
Consequently the distribution of $S$ is not regular.
\end{theorem}

The construction is intentionally elementary.  The marginal survival function
is chosen so that its inverse hazard ratio is $a+x$ on the first interval and
$1-x$ on the second interval.  Thus the marginal virtual value is constant and
then increasing, with no downward jump at the breakpoint.  For the sum of two
copies, we examine only the one-sided neighborhood $s=1+t$ of the interior
point $1$.  In this range the convolution density and survival probability have
closed forms, and differentiating their quotient gives the negative derivative
displayed in \Cref{thm:main}.

The rest of the paper is organized as follows.  \Cref{sec:preliminaries}
records the notation and basic convolution identities used throughout.
\Cref{sec:overview} gives the short proof roadmap and explains why the chosen
parameters make the computation local.  \Cref{sec:main-proof} proves the
regularity of the marginal, derives the local convolution formulas, and
establishes the negative virtual-value slope for the sum.

\section{Preliminaries}
\label{sec:preliminaries}

All random variables below are nonnegative and absolutely continuous.  For a
random variable $X$ with distribution function $F$ and density $f$, write
\[
\overline F(x)=1-F(x)
\]
for its survival function.  We extend densities by zero outside their supports.
The inverse hazard ratio and virtual value are
\begin{equation}
\label{eq:virtual-value}
r_F(x)=\frac{\overline F(x)}{f(x)},
\qquad
\varphi_F(x)=x-r_F(x),
\end{equation}
at points where $f(x)>0$ and $\overline F(x)>0$.  A distribution is regular if
$\varphi_F$ is nondecreasing on this set.  For the compactly supported
piecewise smooth distributions used here, the right endpoint is excluded
because the survival probability is zero, and matching one-sided values at a
breakpoint are read as the value of the displayed formula there.

If $X_1,X_2$ are independent copies of $X$ and $S=X_1+X_2$, write $H$ for the
distribution function of $S$, $g$ for its density, and
$\overline H=1-H$ for its survival function.  The only convolution identities
used below are
\begin{equation}
\label{eq:convolution-identities}
g(s)=\int_{\mathbb R} f(x)f(s-x)\,dx,
\qquad
\overline H(s)=\int_{\mathbb R} f(x)\overline F(s-x)\,dx.
\end{equation}
The virtual value of the sum is
\begin{equation}
\label{eq:sum-virtual-value}
\Phi(s)=s-\frac{\overline H(s)}{g(s)}
\end{equation}
where $g(s)>0$ and $\overline H(s)>0$.  If $\Phi$ has a negative right
derivative at an interior support point, then $\Phi$ is not nondecreasing.  All
logarithms are natural.

\section{Overview}
\label{sec:overview}

The proof gives a direct i.i.d. counterexample.  Rather than trying to control
regularity under an arbitrary convolution, we choose a marginal whose inverse
hazard ratio is completely explicit on two intervals.  The constants
\[
a=\frac1{50},\qquad c=\frac{49}{100},\qquad d=1-c=\frac{51}{100},
\qquad B=\frac{a}{d^2}
\]
are chosen so that $a+c=d$.  This makes the survival function
\[
\overline F(x)=
\begin{cases}
\dfrac{a}{a+x}, & 0\leq x\leq c,\\[6pt]
B(1-x), & c\leq x\leq 1
\end{cases}
\]
fit together continuously, and in fact makes the density pieces agree at the
breakpoint.  The motivation for this normalization is that the inverse hazard
ratio is then $a+x$ on the first piece and $1-x$ on the second piece.  Thus the
marginal virtual value is $-a$ followed by $2x-1$, and these two values agree
at $x=c$.  This proves the marginal regularity in
\Cref{lem:marginal-regular} without any hidden smoothing argument.

The nontrivial part is that convolution regularity depends on the ratio
$\overline H(s)/g(s)$ for the sum, not on the two marginal virtual values
separately.  The proof therefore looks only at a one-sided neighborhood of the
interior point $s=1$ in the support of $S=X_1+X_2$.  Writing $s=1+t$, the
choice $c<1/2$ ensures that for $0\leq t\leq c$ the convolution integrals have
no low-low region.  They split only into low-high, high-high, and high-low
pieces, which leaves elementary formulas for
\[
G(t)=g(1+t)
=
2B\left(\frac{a}{a+t}-\frac{a}{d}\right)+B^2(a+t)
\]
and
\[
T(t)=\overline H(1+t)
=
Ba\left(2\log\frac{d}{a+t}+\frac{a+t}{d}-1\right)
+\frac{B^2}{2}\left(d^2-(c-t)^2\right).
\]
\Cref{lem:convolution-local-formulas} derives these identities directly from
the convolution formulas for the density and survival probability.

The final step is to evaluate the virtual value of the sum,
\[
\Phi(s)=s-\frac{\overline H(s)}{g(s)}, \qquad
\Psi(t)=\Phi(1+t)=1+t-\frac{T(t)}{G(t)}.
\]
Differentiating the displayed local formulas at $t=0$ gives
\[
\Psi'(0+)
=
\frac{31255005}{12510002}
-\frac{6759999}{6255001}\log\frac{51}{2}.
\]
The proof then supplies an elementary sign certificate: $\log(51/2)>3$, so the
last quantity is negative.  Hence $\Phi$ decreases immediately to the right of
the interior point $1$, which is exactly the obstruction to regularity needed
in \Cref{lem:sum-has-negative-slope}.  Combining this with the regularity of
the common marginal proves \Cref{thm:main}: regularity is not closed under
convolution, even for two i.i.d. nonnegative continuous summands.

\section{Proof of the Main Theorem}
\label{sec:main-proof}

Set
\[
a=\frac{1}{50},\qquad
c=\frac{49}{100},\qquad
d=1-c=\frac{51}{100},\qquad
B=\frac{a}{d^2}=\frac{200}{2601}.
\]
The identity $a+c=d$ will make the two pieces of the survival function meet
continuously at $c$.

The first lemma verifies that the prescribed marginal is a regular absolutely
continuous distribution.  The point of the construction is that the inverse hazard ratio
has two elementary pieces whose associated virtual values join without a
downward jump.

\begin{lemma}[Regular marginal construction]
\label{lem:marginal-regular}
Let $F$ be the distribution on $[0,1]$ with survival function
\begin{equation}
\label{eq:marginal-survival}
\overline F(x)=
\begin{cases}
\dfrac{a}{a+x}, & 0 \leq x \leq c,\\[6pt]
B(1-x), & c \leq x \leq 1
\end{cases}
\end{equation}
with $\overline F(x)=1$ for $x<0$ and $\overline F(x)=0$ for $x>1$.  Then
$F$ is absolutely continuous with density
\begin{equation}
\label{eq:marginal-density}
f(x)=
\begin{cases}
\dfrac{a}{(a+x)^2}, & 0 \leq x<c,\\[6pt]
B, & c<x<1,
\end{cases}
\end{equation}
and $F$ is regular.
\end{lemma}

\begin{proof}
Since $a+c=d$, the two displayed formulas for $\overline F$ agree at $c$:
\[
\frac{a}{a+c}=\frac{a}{d}=Bd.
\]
Also $\overline F(0)=1$ and $\overline F(1)=0$.  The function $\overline F$ is
nonincreasing and absolutely continuous on $[0,1]$, so it defines an
absolutely continuous distribution with density
$f(x)=-\overline F'(x)$ at every non-breakpoint point.  This gives
\Cref{eq:marginal-density}.  The two density pieces also
agree at the breakpoint because
\[
\frac{a}{(a+c)^2}=\frac{a}{d^2}=B.
\]

For $0\leq x<c$,
\[
\frac{\overline F(x)}{f(x)}
=
\frac{a/(a+x)}{a/(a+x)^2}
=a+x,
\]
so by \Cref{eq:virtual-value} the virtual value is
\[
x-\frac{\overline F(x)}{f(x)}=-a.
\]
For $c<x<1$,
\[
\frac{\overline F(x)}{f(x)}
=
\frac{B(1-x)}{B}
=1-x,
\]
so by \Cref{eq:virtual-value} the virtual value is
\[
x-\frac{\overline F(x)}{f(x)}=2x-1.
\]
At the joining point, $2c-1=-1/50=-a$.  Hence the virtual value is constant on
$[0,c]$ and increasing on $[c,1)$, and $F$ is regular.
\end{proof}

It remains to show that regularity fails after adding two independent copies.
Only a one-sided neighborhood of the interior point $1$ is needed.  In that
neighborhood the low-low integration region is absent because $1>2c$, leaving
only low-high, high-high, and high-low contributions.

\begin{lemma}[Local convolution formula]
\label{lem:convolution-local-formulas}
Let $X_1,X_2$ be independent with the distribution in
\Cref{lem:marginal-regular}, and let $S=X_1+X_2$.  Let $g$ be the density of
$S$ and let $\overline H(s)=\Pr[S\geq s]$.  For every $0\leq t\leq c$, with
$s=1+t$,
\begin{equation}
\label{eq:local-density}
g(1+t)
=
2B\left(\frac{a}{a+t}-\frac{a}{d}\right)+B^2(a+t)
\end{equation}
and
\begin{equation}
\label{eq:local-survival}
\overline H(1+t)
=
Ba\left(2\log\frac{d}{a+t}+\frac{a+t}{d}-1\right)
+\frac{B^2}{2}\left(d^2-(c-t)^2\right).
\end{equation}
\end{lemma}

\begin{proof}
For $s=1+t$ with $0\leq t\leq c$, \Cref{eq:convolution-identities} gives
\[
g(s)=\int_{s-1}^{1} f(x)f(s-x)\,dx.
\]
Because $s>2c$, there is no interval on which both $x\leq c$ and
$s-x\leq c$.  Splitting at $x=c$ and $x=s-c=d+t$ gives
\[
g(s)=
\int_t^c \frac{a}{(a+x)^2}B\,dx
+\int_c^{d+t} B^2\,dx
+\int_{d+t}^1 B\frac{a}{(a+s-x)^2}\,dx.
\]
The first and third integrals are equal by the substitution $y=s-x$.  Since
$s-2c=a+t$ and $a+c=d$,
\[
g(1+t)
=2B\int_t^c \frac{a}{(a+x)^2}\,dx+B^2(a+t)
=
2B\left(\frac{a}{a+t}-\frac{a}{d}\right)+B^2(a+t).
\]
This proves \Cref{eq:local-density}.

For the survival probability, \Cref{eq:convolution-identities} gives
\[
\overline H(s)=\int_{s-1}^{1} f(x)\overline F(s-x)\,dx,
\]
which holds for $1\leq s\leq 2$.  The same split gives
\[
\overline H(s)
=
\int_t^c \frac{a}{(a+x)^2}B(1-s+x)\,dx
+\int_c^{d+t} B^2(1-s+x)\,dx
+\int_{d+t}^1 B\frac{a}{a+s-x}\,dx.
\]
Since $s=1+t$, the first integral is
\[
B\int_t^c \frac{a(x-t)}{(a+x)^2}\,dx
=
Ba\left(\log\frac{d}{a+t}+\frac{a+t}{d}-1\right).
\]
The middle integral is
\[
B^2\int_c^{d+t}(x-t)\,dx
=
\frac{B^2}{2}\left(d^2-(c-t)^2\right),
\]
and the last integral is
\[
Ba\int_{d+t}^1\frac{dx}{a+1+t-x}
=
Ba\log\frac{d}{a+t}.
\]
Adding these three identities proves \Cref{eq:local-survival}.
\end{proof}

The final lemma differentiates the local formula.  The derivative is evaluated
at an interior point of the support of $S$, so a negative right derivative
forces a local decrease of the convolution virtual value.

\begin{lemma}[Negative convolution virtual slope]
\label{lem:sum-has-negative-slope}
For $S=X_1+X_2$ as in \Cref{lem:convolution-local-formulas}, the virtual value
\begin{equation}
\label{eq:convolution-virtual-value}
\Phi(s)=s-\frac{\overline H(s)}{g(s)}
\end{equation}
is not nondecreasing on the support of $S$.
\end{lemma}

\begin{proof}
Write $\Psi(t)=\Phi(1+t)$ for $0\leq t\leq c$.  By
\Cref{eq:local-density,eq:local-survival}, put
\begin{equation}
\label{eq:G-and-T}
\begin{aligned}
G(t)&=g(1+t)
 =
2B\left(\frac{a}{a+t}-\frac{a}{d}\right)+B^2(a+t),\\
T(t)&=\overline H(1+t)
 =
Ba\left(2\log\frac{d}{a+t}+\frac{a+t}{d}-1\right)
+\frac{B^2}{2}\left(d^2-(c-t)^2\right).
\end{aligned}
\end{equation}
Then $\Psi(t)=1+t-T(t)/G(t)$.  At $t=0$,
\[
G(0)=\frac{1000400}{6765201},
\qquad
G'(0)=-\frac{51980000}{6765201},
\]
and
\[
T(0)=\frac{8}{2601}\log\frac{51}{2}-\frac{9596}{6765201},
\qquad
T'(0)=-\frac{1000400}{6765201}.
\]
Differentiating $\Psi(t)=1+t-T(t)/G(t)$ from the right at $0$ gives
\begin{equation}
\label{eq:slope-value}
\Psi'(0+)
=1-\frac{T'(0)G(0)-T(0)G'(0)}{G(0)^2}
=
\frac{31255005}{12510002}
-\frac{6759999}{6255001}\log\frac{51}{2}.
\end{equation}
We next give a numerical sign certificate using only elementary estimates.
Since
\[
e=1+1+\frac{1}{2}+\sum_{n\geq 3}\frac{1}{n!}
\leq
\frac{5}{2}+\sum_{k\geq 0}\frac{1}{6\cdot 4^k}
=\frac{49}{18}
<\frac{11}{4},
\]
where $n!\geq 6\cdot 4^{n-3}$ for $n\geq 3$, we have $e^3<51/2$ and hence
$\log(51/2)>3$.  Therefore
\[
\Psi'(0+)
<
\frac{31255005}{12510002}
-3\frac{6759999}{6255001}
=
-\frac{9304989}{12510002}
<0.
\]
By \Cref{eq:slope-value} and the definition of the right derivative, there is
some sufficiently small
$\epsilon>0$ such that
\[
\Phi(1+\epsilon)=\Psi(\epsilon)<\Psi(0)=\Phi(1).
\]
Thus $\Phi$ is not nondecreasing on the support of $S$.
\end{proof}

\begin{proof}[Proof of \Cref{thm:main}]
Let $X_1$ and $X_2$ be independent random variables with the common distribution
constructed in \Cref{lem:marginal-regular}.  That lemma shows that the common
distribution is regular.  The proof of \Cref{lem:sum-has-negative-slope}
computes the right derivative of the virtual value of $S=X_1+X_2$ at $1$ as
the quantity in \Cref{eq:slope-value} and proves that it is negative.
Therefore the virtual value of $S$ decreases immediately to the right of $1$,
and the distribution of $S$ is not regular.  This proves that regularity is not
closed under convolution, even for two i.i.d. nonnegative continuously
distributed random variables.
\end{proof}

\bibliographystyle{alpha}
\bibliography{paper}

\end{document}
